\documentclass[12pt,a4paper]{article}

% ---------- Basic packages ----------
\usepackage[utf8]{inputenc}
\usepackage[T1]{fontenc}
\usepackage{lmodern}
\usepackage[margin=1in]{geometry}
\usepackage{graphicx}
\usepackage{array}
\usepackage{setspace}
\usepackage{microtype}

\usepackage{tabularx,adjustbox,array}

\usepackage{xstring}
\usepackage{etoolbox}


% \usepackage[hidelinks]{hyperref}
\usepackage{hyperref}
\usepackage{pdfcomment}
\usepackage{cleveref}

%better bibtex
\usepackage[backend=biber,style=authoryear,sorting=nyt,natbib=true]{biblatex}
\addbibresource{references.bib}

\newbibmacro*{citewithtitle:tooltip}{%
  \printnames{labelname}%
  \addcomma\space
  \printfield{year}%
  \addspace\textbar\addspace
  \printfield{title}%
}

% (Author, Year | Title) with link + tooltip
\DeclareCiteCommand{\citewithtitle}
  {\usebibmacro{prenote}}
  {\printtext[bibhyperref]{%
     \pdftooltip{%
       \mkbibparens{%
         \printnames{labelname},\space
         \printfield{year}\addspace\textbar\addspace
         \printfield[citetitle]{title}%
       }%
     }{%
       \usebibmacro{citewithtitle:tooltip}%
     }%
  }}
  {\multicitedelim}
  {\usebibmacro{postnote}}
% Make natbib-like commands available:
\let\citep\parencite
\let\citet\textcite

\makeatletter
\edef\dictfile{macros/keyword.tex}
%\typeout{Using dictionary: \dictfile}%
\input{\dictfile}
\makeatother

% --- Load the dictionary with a safe fallback ---
\makeatletter
\edef\dictfile{macros/dictionary-\DICT.tex}
%\typeout{Using dictionary: \dictfile}%
\input{\dictfile}
\makeatother

% --- Simple conditional helpers based on \DICT ---
\newcommand{\ifdev}[1]{%
  \IfStrEq{\DICT}{development}{#1}{}%
}

\newcommand{\ifsus}[1]{%
  \IfStrEq{\DICT}{sustainability}{#1}{}%
}


% centered, stretchable column
\newcolumntype{Y}{>{\centering\arraybackslash}X}
\newcommand{\logoheight}{1.1cm} % tweak if needed


\setstretch{1.2}
\setlength{\parindent}{0pt}
\setlength{\parskip}{0.6em}

% For even spacing of logos
\newcolumntype{C}[1]{>{\centering\arraybackslash}m{#1}}

% Make article say "Table of Contents"
\renewcommand{\contentsname}{Table of Contents}

% -----------------------------------
% Editable metadata
% -----------------------------------
\newcommand{\ThesisTitle}{\TITLE}
% \newcommand{\ThesisSubtitle}{SUBTITLE}
\newcommand{\ThesisCity}{[City]}
\newcommand{\ThesisDate}{[Date]}

\newcommand{\AuthorName}{Daniel Morgan Rivers}
\newcommand{\AuthorAddressTwo}{Berlin, Germany 10559}
\newcommand{\AuthorEmail}{rivers@uni-potsdam.de}
\newcommand{\MatricNo}{829112}

\newcommand{\FirstReviewer}{Prof. Christian Kuhlicke}
\newcommand{\SecondReviewer}{Dr. Ivan Kuznetzov}
\newcommand{\MonthName}{%
  \ifcase\month
  \or January\or February\or March\or April\or May\or June%
  \or July\or August\or September\or October\or November\or December%
  \fi
}
\newcommand{\DateDDMonthYYYY}{\number\day~\MonthName~\number\year}

\begin{document}
\pagenumbering{gobble}
\thispagestyle{empty}

% ---------- Title page ----------
\begin{center}

    % Desired order: RIFS, UFZ, AWI, PIK, University crest
    \noindent
    \begin{tabularx}{\textwidth}{@{}YYYYY@{}}
    % \adjustbox{max height=\logoheight, max width=\linewidth, keepaspectratio}{\includegraphics{assets/rif.jpg}} &
    \adjustbox{max height=\logoheight*1, max width=\linewidth*1, keepaspectratio}{\includegraphics{assets/helmholtz.png
      }} &
    \adjustbox{max height=\logoheight*1, max width=\linewidth*1, keepaspectratio}{\includegraphics{assets/awi.jpg}} &
    % \adjustbox{max height=\logoheight, max width=\linewidth, keepaspectratio}{\includegraphics{assets/pik.png}} &
    \adjustbox{max height=\logoheight*2, max width=\linewidth*2, keepaspectratio}{\includegraphics{assets/potsdam.png}}\\
    \end{tabularx}


    \vspace{1.8cm}

    
    \begin{center}
        \begin{flushleft}
        University of Potsdam\\
        Faculty of Science\\
        Institute of Environmental Science and Geography\\
        Institute of Physics and Astronomy\\
        \textbf{Climate, Earth, Water, \& Sustainability\\[4em]}
        \end{flushleft}
    
        {\large \textbf{Master Thesis}}\\
        for the award of the academic degree\\
        \textbf{Master of Science (M.sc.)}\\
        at the University of Potsdam\\[3em]
    
        \textbf{\LARGE \ThesisTitle}\\ [3em] %if no subtitle
        %\textbf{\large SUBTITLE}\\[6em]
    
        Potsdam, \DateDDMonthYYYY \\[9em]
    
        \begin{flushleft}
            \textbf{Submitted by:}\\[-0.2em]
            \AuthorName\\
            % \AuthorAddressOne\\
            % \AuthorAddressTwo\\
            \AuthorEmail\\
            Matriculation No.: \MatricNo

            % \vspace{1.0cm}
            First reviewer: \FirstReviewer\\
            Second reviewer: \SecondReviewer
        \end{flushleft}

    %     \textbf{Submitted by:}\\
    %     Louis Lacassagne\\
    % %    [Address 1]\\ unless it is a legal requirment, i dont think adding the adress is useful at all
    % %    [Address 2]\\
    %     lacassagne@uni-potsdam.de\\
    %     Matriculation No.: 820718\\[3em]
    
    %     \textbf{First reviewer:} Prof. Dr. Georg Feulner\\
    %     \textbf{Second reviewer:} Julius Eberhard, M.Sc.
    %     \end{flushleft}
    \end{center}
    \clearpage

    % {\Large University of Potsdam\\
    % Faculty of Science\\
    % Institute of Environmental Science and Geography\\
    % Institute of Physics and Astronomy\\[0.3em]
    % % Keep this upright (no italics) for a cleaner look
    % Climate, Earth, Water, \& Sustainability}

    % \vspace{2.0cm}
    % {\LARGE \textbf{\ThesisTitle}}\\[0.4em]
    % % {\large \ThesisSubtitle}

    % \vspace{2.0cm}
    % {\Large \textbf{Master Thesis}}\\[0.6em]
    % {\large for the award of the academic degree}\\[0.1em]
    % {\large \textbf{Master of Science (M.\,Sc.)}}\\[0.2em]
    % {\large at the University of Potsdam}

    % \vspace{1.8cm}

    % {\large \ThesisCity, \ThesisDate}

    % \vfill

    % \begin{flushleft}
    %     \textbf{Submitted by:}\\[-0.2em]
    %     \AuthorName\\
    %     \AuthorAddressOne\\
    %     \AuthorAddressTwo\\
    %     \AuthorEmail\\
    %     Matriculation No.: \MatricNo

    %     % \vspace{1.0cm}
    %     First reviewer: \FirstReviewer\\
    %     Second reviewer: \SecondReviewer
    % \end{flushleft}
\end{center}
\clearpage
% ---------- Main matter ----------
\section*{Abstract}
\subsection*{Abstract in English}
The field of decision science has made rapid strides with the introduction of techniques of deep learning in natural language processing as a tool for improving accuracy and calibration of economic and geopolitical forecasts, predicting scientific outcomes, and in domains like stock market prediction and diagnosing diseases. Several recent methods and improvements have been made to the predictive ability of large language models (LLMs) calibrated via fine-tuning techniques. Because some of these powerful forecasting techniques have not yet been brought to bear on the important problem of impact prediction \ABSTRACT, I fine-tuned two LLMs (Llama 70B and ChatGPT 3.5) to perform such forecasts, training on thousands of abstracts from the scientific literature. I find that the best-performing system ([Llama 70B / ChatGPT 3.5], averaging X parallel predictions) can forecast XX\% of interventions correctly compared to YY\% using a random forest model baseline, and ZZ\% from published ex-ante predictions from the scientific literature. This work has wide-ranging applications both to improve \DOMAIN\ forecasting, as well as in adjacent areas and improving decision making in policy contexts.

\subsection*{Abstract in German}
Lorem ipsum ...
\clearpage

% ---------- Contents & front matter ----------
\pagenumbering{roman}
\setcounter{page}{1}
\tableofcontents
\clearpage

\pagenumbering{arabic}
\setcounter{page}{1}





















\section{Background: The Science of Forecasting}
\subsection{Introduction}

\textbf{Background} 

The earth system sciences concern the complex interaction between biological, chemical, physical, and anthropogenic processes. A broad goal of the earth system sciences is to model and accurately predict the outcomes of interventions with regard to the environment and its impact on humans. Much of the progress in earth system science has been on linking these complex phenomena into large models, such as integrated assessment models (IAMs), computable general equilibrium models (CGEs), or agent-based models (ABMs). While many attempts have been made to model specific subsystems within the earth system, such as the carbon cycle, environmental and economic linkages, or understanding human impacts in the climate-water-food nexus, there have been few attempts to create a comprehensive model which can predict quantitative or qualitative outcomes of a wide range of cross-domain interventions in the earth system which could be described in natural language.

Thus far, models in the earth system sciences have relied on mechanistic, theoretically-based models of the underlying complex systems. 
However, this is not the only way to predict outcomes - statistical, or ``Machine Learning'' (ML) outcomes, while lacking the rigorous mechanistic underlying processes, have recently been shown to perform better in many cases to model complex systems such as language modelling \citewithtitle{brownLanguageModelsAre2020}, protein folding \citewithtitle{jumperHighlyAccurateProtein2021}, biodiversity protection \citewithtitle{silvestroImprovingBiodiversityProtection2022}, and weather forecasting \citewithtitle{lamLearningSkillfulMediumrange2023}.

The collective failure of the scientific community to model complex outcomes in the earth system has severe implications. For example, work from \citewithtitle{stechemesserClimatePoliciesThat2024} has demonstrated that out of 1500 policies between 1998 and 2022, only 68 had statistically significant causal effect to reduce country emissions with a 99\% or higher confidence. Furthermore, they find that more than four times the effort witnessed so far in emissions reductions from implementing more successful policies in line with past reductions would have to be exerted to close the emissions gap to remain below 2 degrees C in global temperature rise. Broadly, their findings support the claim that even when climate policy is implemented, it is largely ineffective, and in the future it will need to be much more effective to avoid dangerous levels of CO$_2$ concentrations. In terms of biodiversity,  achieving sustainability cannot be met by current trajectories, and goals for 2030 and beyond may only be achieved through transformative changes across economic, social, political and technological factors \citewithtitle{watsonGlobalAssessmentReport2019}. As of 2022 pollution remains responsible for approximately 9 million deaths per year, corresponding to one in six deaths worldwide \citewithtitle{fullerPollutionHealthProgress2022}.

While much scientific effort has been expended on understanding underlying systems, much less effort has been directly focused on predicting which specific policies, if enacted, would realistically improve outcomes on the indicators of interest in the earth system sciences. Meanwhile, examples exist in the literature where regulation can greatly reduce or even eliminate environmental problems - the Montreal protocol has met with great success in closing the hole in the ozone layer [CITE IF KEEP THIS SENTENCE!]. 


IAMs have shown promise in modelling outcomes of specific policies, with the disadvantage that they are harder to use and set up, require a high computational power and expertise to use effectively, and are not rigorously benchmarked on large databases of existing interventions and associated outcomes. For any user-defined policy package (for example, introducing efficient clean-burning cookstoves in India), GAINS can calculate the reduction in emissions (PM-2.5, NO$_x$, CO$_2$, etc), the improvement in ambient air quality, and the health impacts such as lives saved from lower PM-2.5 exposure \citewithtitle{CosteffectiveControlAir2011}. Other IAMs include the MIT Emissions Prediction and Policy Analysis (EPPA) model, which requires manually entering assumptions of the effects of policies into models of the world economy, calculates the implications on health and runs a CGE to estimate the economic effects \citewithtitle{MITEmissionsPrediction}. 
In the domain of biodiversity, an ML-based framework called CAPTAIN uses a reinforcement learning (RL) agent coupled with a spatially explicit ecosystem simulation to statistically learn which areas to protect over time in order to maximize species survival under budget constraints, to maximize cost-effectiveness in protecting biodiversity \citewithtitle{silvestroImprovingBiodiversityProtection2022}. 
Other techniques used to predict outcomes of interventions include linear optimisation combined with econometric theory, such as the Open Source Energy Modelling System (OSeMOSYS). OSeMOSYS simulates energy production and consumption under policy constraints including a model of the energy grid. By incorporating physical and known constraints, such models have the potential to predict outcomes of policy interventions over longer time horizons \citewithtitle{OSeMOSYSOpenSource2011}. An even more fine-grained, bottom up approach of modelling intervention outcomes is possible. For example, combining bio-economic farm optimization models with ABMs, researchers have modelled evolution of pesticide-related risks for the country of Switzerland \citewithtitle{dueriModelingImplicationsPolicy2024}.

Despite all these examples, in many relevant sub-domains, such as climate policy, ex ante analysis of mitigation action and of mitigation plans is limited \citewithtitle{intergovernmentalpanelonclimatechangeipccMitigationDevelopmentPathways2023}. Given the overwhelming complexity of the earth system, and the corresponding failures to properly model many of the system components in the earth system and especially how they interact with human interventions, complementing mechanistic understanding and prediction with ML approaches is urgently needed. 


\textbf{Proposal}

This thesis proposes the use of LLMs to close this research gap. I will briefly review current progress in specific domains of outcome prediction in the earth system sciences, and then discuss progress with LLMs in adjacent domains. To my knowledge, there has been no attempt at predicting real-world outcomes of interventions in the earth system sciences while also rigorously quantifying the skill of such a system.

Within the domain of LLM use, there has been some progress. A recent tool called ``climsight'' summarizes and aggregates information about climate adaptation and mitigation \citewithtitle{koldunovLocalClimateServices2024}, but stops short of making predictions towards adaptation. Machine learning and LLMs have been used to collect over 80,000 articles about climate adaptation and provide analysis about which areas of implementation are lacking and point out gaps in attention towards promising categories of policies.

Finally, limited work has been done using LLMs such as ChatGPT-4 to serve as data sources for policy deliberation and multi‑criteria assessment of climate and sustainability interventions, finding GPT-4 is in rough agreement with the policy rankings of human experts for the expected outcomes \citewithtitle{binaLargeLanguageModels2025}. However, very little is done to improve on GPT4's abilities, the assessment was made on only a few dozen generic policy examples, and no attempt was made to compare outcomes between these policies and real-world outcomes. Despite these limitations, the findings are promising. For multiple criteria decision making (MCDM), GPT-4 provided a useful collaborative starting point, eased the process of considering multiple criteria effectively, and aided policy deliberation on climate change and sustainability.


\subsection{Prediction Markets and Superforecasting}

In recent years, significant progress has been made on accurate near-term forecasting outside of specific domains. The most promising approaches appear to be a mix of prediction markets, and specialized, trained experts known as ``superforecasters'' \citewithtitle{tetlockSuperforecastingArtScience2015}. Prediction markets have gained recent prominence in the domain of geopolitical forecasting, with significant volumes of transactions on predicting future geopolitical outcomes with a broad purvue, including election results, the outcomes of treaties, or whether a regime will topple. Predictions are typically more accurate than a guess of 50\% chance on a given outcome only within approximately one year time horizon, with the accuracy notably improving as the event reaches only a few months until the question is closed \citewithtitle{tetlockSuperforecastingArtScience2015} NOTE: CHECK+REFINE THIS STATEMENT.  In a broad range of complex, human-involved outcomes, prediction markets are superior to expert analysis. In the words of the economist Robin Hanson, ``racetrack market odds improve on the prediction of racetrack  experts; orange juice commodity futures improve on government weather forecasts;  stocks fingered the guilty firm in the Challenger crash long before the official NASA  panel; Oscar markets beat columnist forecasts; gas demand markets beat gas  demand experts; betting markets beat Hewlett Packard official printer sale  forecasts; and betting markets beat Eli Lily official drug trial forecasts.''  \citewithtitle{hansonShallWeVote2013}. 

However, prediction markets have demonstrated that a smaller subset of forecasters in the market, known as ``superforecasters'', are statistically much better forecasters than the prediction market, and ensembling these forecasters and letting them exchange information among themselves leads to higher accuracy predictions than prediction markets alone \citewithtitle{mellersIdentifyingCultivatingSuperforecasters2015}. Due to the relatively high expense and human effort required to organize superforecasters, the efforts have been largely focused on specific geopolitical and economic questions, some of which may fall under the domain of intervention impact in the earth system sciences, although most point to broad trends where information may be gathered from the news and informal internet searches, and deep expertise in any single domain would not be required for an accurate forecast. In fact, in terms of ``calibration'', superforecasters usually beat domain experts in their own fields by maintaining a broad sense of good judgement and cultivating a trained skill at accurately estimating prediction probabilities, rather than overly relying on a single strategy (such as econometric analysis, or specific statistical methods).

\subsection{LLM Forecasting of Real-World Outcomes}

Assembling superforecasters to predict the efficacy of interventions such as laws, specific interventions such as introduction of cleaner burning ovens, or regulations on air quality would be prohibitively costly. Instead, this work focuses on the mimicking of techniques known to be effective for superforecasters with LLMs. 

As might be expected given the lack of world knowledge and complex reasoning abilities of LLMs, simply replacing a crowd of humans with a crowd of untrained LLMs does not generally outperform the crowd average \citewithtitle{abolghasemiHumansVsLarge2025} \citewithtitle{schoeneggerLargeLanguageModel2023} CHECK THESE REFS.

While the base models appear to underperform compared to crowds of humans, more recent work on the question has shown that increasing model reasoning ability increases the forecasting accuracy \citewithtitle{yanq.serajr.hej.mengl.andsylvaint.AutoCastEnhancingWorld2024}, and that with proper techniques and careful prompting, LLMs will approach or sometimes exceed accuracy of superforecasters on questions with a high degree of context and with proper ensembling and fine-tuning of the LLM system \citewithtitle{halawiApproachingHumanLevelForecasting2024}. \citewithtitle{WisdomSiliconCrowd} \citewithtitle{abolghasemiHumansVsLarge2025} \citewithtitle{yanq.serajr.hej.mengl.andsylvaint.AutoCastEnhancingWorld2024}.  NOTE: CHECK THIS STATEMENT

The best results are achieved by capitalizing on the broad world-knowledge of LLMs and the augmentation of their knowledge in high-news or near-term contexts \citewithtitle{halawiApproachingHumanLevelForecasting2024}. A few crucial improvements to the base-level prediction ability can be applied to approach superforecasting level calibration and accuracy. These improvements are: 
\begin{enumerate}
\item Fine-tuning the LLMs to replicate the format of good forcasts, using hundreds or thousands of correct forecasts as the fine-tuning dataset
\item Have the LLM integrate relevant and timely information into the context to improve the forecast (such as NewsCatcher and Google News API)
\item Have the LLM split questions into sub-questions before being used to query news API
\item Prompting techniques: Have the LLM think step-by-step, rephrase the question to improve comprehension, and reason over chains of crafted prompts to ensure sufficient reasoning effort has gone into the answer
\item Reduce error rates by aggregating the final predictions (``Wisdom of the crowd'')
\end{enumerate}

In one similar work, the technique of Chain of Thought (CoT) has been used to improve the reasoning abilities of GPT4 in predicting the outcome of 1261 conclusions from 276 papers which analyze the real-world outcomes of field experiments in the social sciences. While not specifically investigating outcomes with relevance in the earth system sciences, they do investigate the prediction ability for the impact of educational incentives, household finance behavior, healthcare enrollment, and financial planning. Remarkably, over the 1261 outcomes, 78\% were predicted accurately by the system \citewithtitle{chenPredictingFieldExperiments2025}.

%In another work, it was shown that LLMs can mimic the respondents of surveys in the social sciences, besting expert forecasts of expected outcomes in the domain {hewittPredictingResultsSocial}.  %NOTE: THIS SEEMS INSUFFICIENTLY RELEVANT

Lastly, a recent study found that LLMs can correctly predict outcomes in scientific domains such as predicting results of papers in neuroscience, not just in geopolitical or economic forecast contexts \citewithtitle{xLargeLanguageModels2025}. Notably, this result used the raw probabilities generated by the language model rather than explicit reasoning, and for this reason was able to use very small language models compared to GPT4 as was used in most other studies. Because language models work by assigning a probability of each word,  multiplying the probabilities of all the words multiplied in the entire abstract allows researchers to compare the multiplied probability of the real abstract to the multiplied probability of the fabricated abstract directly, without having the language model generating any text involving reasoning or chain of thought. 

Another study found a similar result with regards to publications in the domain of AI algorithms, finding their system beating human experts in predicting the ability of an AI algorithm to improve on the state of the art performance in AI models \citewithtitle{wenPredictingEmpiricalAI2025}. In this domain, the researchers use a sophisticated framework with RAG and fine-tuning, and still find a good result.

While much cheaper than prediction markets or IAMs, LLMs are also more expense than simpler ML models. When attempting to forecast whether results and effect sizes replicate in social sciences, simple neural network classifiers trained on millions of scientific abtracts and hundreds of full texts, the unordered semantic vectors of the words in the abstracts of the papers, combined with statistical  were sufficient to approach prediction market level accuracy of approximately 70\% accuracy in predicting which paper results would replicate, despite lacking fundamental logical relationships between words in the text or any deeper language comprehension of the methods of the abstracts \citewithtitle{yangEstimatingDeepReplicability2020}. This finding mirrors that of the neuroscience study \citewithtitle{xLargeLanguageModels2025} which finds that reasoning  ability is not strictly to predict the outcomes in neuroscience abstracts. 

Insofar as identifying whether results from social science papers will replicate is a similar task as forecasting the impact of an intervention in the earth system sciences, we can be encouraged that statistical and categorical aspects of the interventions should be sufficient to identify the likely success of real-world outcomes, and remain skeptical that LLMs are strictly necessary to rival humans at predicting categorical outcomes, where ML may be sufficient. However, ML models certainly have a lower upper bound in potential accuracy than a full reasoning model, and computational resources are not so restricted that LLMs could not be used. Furthermore, ML models using simple semantic vectors cannot produce free-form predictions of outcomes like LLMs, limiting the flexibility of their application in real-world use-cases.

If LLMs are able to approach or surpass human ability in predicting unpublished results in complex domains of predicting which techniques in improving state of the art AI system performance, predicting the outcomes of neuroscience papers, predicting social science replicability, or predicting geopolitical events such as election results, then it stands to reason that they may be able to predict the outcomes of interventions in the earth system sciences. While geopolitical forecasting may not be amenable to scientific techniques, neuroscience and AI algorithm improvements certainly are - yet LLMs still beat human experts in these domains. Furthermore, LLM systems are far simpler to use, and far less costly to run and maintain than IAMs, CGEs, or ABMs, while having the benefit of human-interpretable reasoning and the ability to be extremely flexible as to their domain of application. Finally, given their low cost to use, LLMs can often be used as starting points or augmentation to expert judgement in ex-ante outcome prediction, rather than being the sole source of judgement about expected intervention outcomes, and the collaboration often bests forecast accuracy when compared to expert forecasts or LLM forecasts alone  \citewithtitle{schoeneggerAIAugmentedPredictionsLLM2025}\citewithtitle{schoeneggerPromptEngineeringLarge2025}.

The use of LLMs to inform decision making for promising interventions in the earth system comes with several downsides. Notably, LLMs do not reason like humans, and are prone to ``hallucinations''  where facts are fabricated. These hallucinations can be either factual fabrications attributed to external source material, or false statements which come intrinsically from the model \citewithtitle{huangSurveyHallucinationLarge2025}. As mentioned previously, LLMs are more computationally costly than other machine learning methods. There are also issues (which we will leave for the \hyperref[sec:conclusion_outlook]{Conclusion \& Outlook} section) with overly trusting LLMs, false beliefs from users of LLMs that they are less biased than humans or not biased at all, and issues with AI safety, if LLMs begin to replace or distort, rather than augment, human decision making.

% \subsection{Predicting Outcomes with AI \DOMAINCAPITALIZEDINTERVENTION} % DON'T THINK I'LL KEEP THIS SECTION
\section{Methods for LLM Forecasting}
\subsection{Data Sources}
\subsection{Selecting LLM's for Forecasting \DOMAINCAPITALIZEDINTERVENTIONTWO }
\subsection{Baseline Measures to Compare Against LLM Forecasts}
\subsection{Methods for LLM Forecasting}
\subsection{Grading Forecast Accuracy}

\ifdev{ \subsection{NOTE: YES THEY REALLY WANT THIS: Predicting the Confidence of a Forecast}}
\section{Results \& Discussion}
\ifdev{\subsection{Database of evaluations}}
\subsection{Strengths and Weaknesses of This Forecasting System}
\subsection{Evaluation of Techniques for Improving Forecast Accuracy}
\subsection{The Risk of Trusting This Forecasting System}
\clearpage
\section{Conclusion \& Outlook}
\label{sec:conclusion_outlook}

\subsection{The State of AI and LLMs}
\subsection{The Promise and Capabilities of AI Forecasting}

As a clear disclaimer: **LLM's are not in general superior to humans at forecasting as of May 2025.** At the same time, their forecasting ability for short-term predictions is closing in at a rapid pace as AI capabilities have advanced [source]. Furthermore, predictions with a significant number of relevant news articles or very near to the date of a forecasting resolution can best teams of trained forecaster's aggregate predictions in prediction accuracy \citewithtitle{halawiApproachingHumanLevelForecasting2024}. It is currently unknown to what extent the ability of AI systems to forecast geopolitical and economic events can be extended to forecasting the impact of interventions with implications in the earth system sciences. Exploring this domain opens a promising avenue to improve the efficacy of interventions in the earth system sciences. In the remaining section, we discuss the beneficial aspects of the system developed, as well as the potential dangers or risks this system may pose.

One co-benefit of a system fine-tuned on earth system sciences is that by its cross-domain nature, the LLM will be able to identify a wide range of likely outcomes, and the degree of effect of those outcomes, on a wide range of quantitative and qualitative outcomes. When implementing interventions, researchers, policy-makers, and decision makers must always consider many relevant outcomes of their interventions. The similarity and vector search of the system allow users to quickly identify relevant documentation as well as outcomes of similar scientific research most relevant to their proposed intervention.

Another benefit of the system is that AI's typically excel in domains where human experts are particularly challenged: when there is a very large range of relevant data or when predictions about the effect of an intervention involve carefully calibrated probabilities. AI's can also perform predictions in a way that human experts can learn from: introducing one piece of information can be used to quantify the effect on AI forecasts. AI forecasts can be ensembled arbitrary and at relatively little expense compared to humans. 

\subsection{Risks, Biases, and Limitations}
However, there are clear risks of using AI for evaluating the likely outcomes of interventions in the earth system sciences. The most obvious issue may be that while AI can be accurate in some domains, current AI systems do not accurately present their confidence in their answers and can completely hallucinate events and facts which have no grounding in reality. The result is a misleading analysis, which in the space of earth system sciences may lead to significant risks. Policy makers may trust AI more than is justified by its performance, or view it as an unbiased source, despite nearly all current AI systems having a well-documented political bias acknowledged by both the political left and political right [source]. 

Another risk is that scientists may not perform research deemed to be unlikely to succeed, and thus the range of explored outcomes may be narrowed to the outcomes known to work in the past or deemed to be likely to be successful by the AI system. 

While AI may be able to calibrate itself on many different domains and automatically pull in relevant information, it currently lacks the ability to reliably perform complex mathematical calculations or run long-term analysis. Furthermore, as AI becomes more advanced there is significant concern in the technology community that it may form its own goals and intrinsic values, out of alignment with its human operators. An AI that advises on AI policy may in fact present a conflict of interest, even if the AI is simply using heuristics mimicking human tendencies towards self-preservation and in-group preferences. 

Finally, without the full text, there is a risk that the policy forecasting aspect may be quite limited. Without a sense of the scope of an intervention, which would not reliably be indicated in the abstract, the degree of impact of an intervention may difficult to ascertain by any forecasting system.

\subsection{System Design and Risk Mitigation}
We address these concerns by noting that as AI begins to become more accurate and lower cost than human researchers at forecasting the impact of policy outcomes, it becomes ever more important to have specifically designed systems that take steps to reduce the dangers of AI systems. We believe the system developed clearly fulfills this criterion. The system we use in this work specifically provides credible, peer-reviewed scientific information and news from reputable sources to the AI, rather than relying on general internet search as many current AI providers rely on.   [OPTIONAL: Furthermore design our system to be calibrated via fine-tuning, meaning that some of the reliability concerns may be ameliorated. ] As AI systems advance, there appears to be a progression towards more agentic systems with more clear intermediate goals. A misalignment with human preferences (an example in this work might be downplaying the CO2 effects of building more AI systems in order to increase the number of AI systems as an in-group preference) may occur and be missed by humans with extremely long thought chains and insufficient detection of misalignment. Our system by contrast allows the user to inspect the series of logical deductions performed by the model and view available sources the model used as scientific reference material. The system has been specifically quantified in terms of its bias, allowing users to have full knowledge of the likely failure modes when using the system, often absent in generally available AI chat interfaces. [ OPTIONAL: with an explicit attempt to correct these biases via fine-tuning, syncophantic behavior is also reduced compared to RLHF models.].
Another risk is that papers tend to have a bias, and the model will learn to replicate that bias. Papers are much more likely to have "significant" results than mixed effect or no effect. The optimistic bias towards positive bias published in journals should mean we interpret the prediction of the model cautiously, with knowledge that it will likely present a more optimistic version of the outcomes than is justified from a neutral observer's perspective. In order to counteract this risk, we are also looking at the accuracy of the quantitative result of the intervention, which is more valid to compare between abstracts and has a relatively smaller publisher bias [source]. 
Finally, much of the promise of the AI forecasting approach relies on models continuing to become lower cost and more performant in general domains. While multiple empirical trends and the longstanding success of Moore's law clearly indicate this should continue, it is by no means guaranteed. If AI models cease to improve on relevant metrics, or otherwise become increasingly biased or unreliable, much of the promist of an AI forecasting tool for estimating interventions in the earth system sciences goes away. Despite this risk, the system remains useful and informative for the scientific and public policy community as it provides a system with sources proven to provide useful information for the evaluation of policy outcomes, and introduces a framework by which the impact of interventions can be broken down for more accurate predictions. While there is a possibility that AI's may never reach the capabilities of humans in integrating the disparate sources of information, automated information search and a new tool that can synthesize relevant information can be a powerful tool for scientists and policy makers.
Forecasting has the distinct benefit of disallowing training on any particular benchmarks and is a rather difficult-to-game metric compared to standard LLM performance benchmarks. It becomes increasingly useful to society to understand what the true capabilities of LLM's are and the rate of their improvement, both for the regulation of dangerous AI capabilities and the improved understanding where AI may be capable enough for reliable use in various critical domains such as automated medicine and driverless vehicles.

\subsection{Broader Applications and Vision}
The codebase and research done here can be repurposed from specifically earth system science, to other domains where impact forecasts are clearly useful. A similar system with an expanded set of abstracts and data could be used with relatively little modification in domains such as public health, financial policy, and in a more general way to provide predictions for scientists about likely qualitative and quantitative outcomes of their scientific studies. The success of the model demonstrates that a great deal of opportunity to synthesize scientific findings and improve decision making on an institutional level is policy.
One particularly promising avenue for expansion of the system would be as an application to Futarchy first proposed by Robin Hanson. Futarchy proposes to use prediction markets to allow policy makers or the general public to only have to agree on what they value and quantify as utility, not on how to maximize that utility. Several prediction markets in parallel are formed, creating a zero-sum game financially rewarding players that best predict the utility outcome conditional on a policy being implemented. To the extent that complex public policy can ever be reduced to a single utility function, that this function can be agreed on by a quorum of policy makers, Futarchy could significantly reduce gridlock and polarization in politics, at least in the domains in which the necessary conditions are useful and possible. In essence, Futarchy aids policy makers in coming to agreement on how to implement policies by reducing the scope of disagreement to what the set of possible policy implementations could be and how they would choose to quantify a successful outcome.
If and when the system proposed is shown to exceed human ability in predicting policy, or if it can be shown that the system can be complementary to human predictions, cheaply improving their accuracy, this system could be integrated to a scheme for futarchy by replacing or augmenting prediction markets. This may be especially helpful in use-cases where AI succeeds and prediction markets fail: very low probabilities over long time periods (as the winners may choose to invest their money on a higher-return investments), predictions about long-run outcomes that are difficult to gain information about, particularly contentious outcomes, or issues where markets may be biased by particularly wealthy individuals who come in very late in the market and buy many more shares than expected.

\subsection{AI Scientist-like idea}
Extending the system for searching for high-impact policies is possible, rather than simply using the fine-tuned model for forecasting. While use of reasoning models outside of the domain in which they are trained for often reduces their performance, it still may be possible to re-train the model for these use cases. For instance, the model could be prompted to generate many policy options for a given country to reduce CO2 emissions, and each idea could have the emissions reduction forecasted. Seeding the model with many similar policies and suggesting that it think of a wide range of options may allow for consideration of a wide range of policy options. Next, only the ideas which are forecasted to have high emissions would be suggested to the user of the system. Such a system would be similar to the ``AI Scientist'' released by Google which iteratively generates new hypotheses and reasons over the hypotheses to discover better scientific theories behind biological phenomena \citewithtitle{luAIScientistFully2024}.

\subsection{Ideation: Extensions and other applications}
\begin{itemize}
    \item Improving prospects of futarchy to improve governance
    \item Understanding how different sources of information contribute to effective forecasting of impact
    \item Before the forecasting at all: collecting the information for forecasting all in one place, both resources to make reasonable forecasts, as well as creating structure out of unstructured papers in earth systems science
    \item Creating general hierarchies of impact for different categories of interventions
    \item Ability to create "unbiased" forecasts that are both evidence based and listened to by both sides of the political spectrum
    \item Increasing democratic understanding of the likely effects of laws from third party sources: allows non-experts to assess the efficacy of elected officials in accomplishing their goals
    \item Automated scoring of introduced legislation 
    \item Sufficient statistics to introduce confidence bars on the effects of political outcomes
    \item Leveraging the advance of AI for good
    \item Constraining the use of AI in a scientifically valid, constrained manner, which minimizes the risk that AI biases themselves influence policy decisions.  
    \item Automated feedback on proposed interventions (registered studies): what are the likely things this has impact on? What are some relevant papers for their proposal?
\end{itemize}






























\clearpage
\section*{Works Cited}
\printbibliography

% \addcontentsline{toc}{section}{Works Cited}
% Add your references or switch to biblatex later.

\clearpage
\section*{Declaration of Academic Integrity}
\addcontentsline{toc}{section}{Declaration of Academic Integrity}
Ich, Morgan Rivers, erkläre hiermit, dass diese Arbeit das Ergebnis meiner eigenen Arbeit ist. Ich danke für die Unterstützung, die ich bei der Erstellung dieser Arbeit erhalten habe, und für die verwendeten Quellen.
\emph{(I, Morgan Rivers, hereby declare that this thesis is the product of my own work. All the assistance received in preparing this thesis and the sources used have been acknowledged.)} \\
% Example (English): I hereby declare that this thesis is the product of my own work.
% When uploading your thesis in PULS, the included Statement of Originality does not necessarily have to be signed.
% \newpage
\\
Potsdam, \DateDDMonthYYYY

\end{document}